\documentclass[10pt,a4paper]{amsart}
\usepackage[margin=19mm]{geometry}
\usepackage{amsmath,amssymb,mathtools}
\usepackage[scr=rsfs,cal=euler]{mathalfa}
\usepackage{xfrac}
\usepackage[unicode,hidelinks,bookmarks=false]{hyperref}
\newcommand{\ie}{i.e.\ }
\newcommand{\cf}{cf.\ }
\newcommand{\resp}{respectively\ }
\newcommand{\Subsection}{\subsection}
\providecommand{\nobreakdash}{\nobreak\hbox{-}}
\setlength{\emergencystretch}{2em}
\multlinegap0pt
\usepackage{xcolor}
\usepackage{tikz}
\usepackage{enumerate}
\usepackage[shortlabels]{enumitem}
\usepackage{mdframed}
\usepackage[normalem]{ulem}
\usepackage{amssymb}
\usepackage{cancel}
\usepackage{dsfont}
\usepackage{braket}
\usepackage{mathtools}
\newtheorem{prop}{Proposition}
\newcommand\brak[1]{\big \langle #1 \big\rangle}
\newcommand{\gr}[1]{{\color{gray} #1}}
\title{Enumeration of general planar hypermaps with an alternating boundary}
\author{Valentin Baillard, Ariane Carrance, Bertrand Eynard}
\date{Source: arXiv:2603.28571v1 (CC BY 4.0)}
\begin{document}
\maketitle

\noindent\emph{Source and licence.} Enumeration of general planar hypermaps with an alternating boundary. Version arXiv:2603.28571v1 (CC BY 4.0), distributed by arXiv under the Creative Commons Attribution 4.0 International licence, http://creativecommons.org/licenses/by/4.0/, as recorded in the arXiv licence metadata for this exact version and shown on the abstract page https://arxiv.org/abs/2603.28571v1. Full licence notice: \texttt{licenses/arxiv-2603.28571-CC-BY-4.0.txt}.\par\smallskip

\noindent\textit{Editorial scope.} Counted unit: the article's prop eq:first-equation-in-M (source0.tex lines 574--584). The article's complete original proof of it is delivered with it (source0.tex lines 586--643, one \texttt{proof} environment of 58 lines). No mathematics is written, repaired or re-derived: the statement and the proof are the article's own lines, in the article's own notation, and no retained line is substituted or re-flowed.  No further source-local context is needed: every cross-reference of every retained line resolves inside the two delivered spans.  Nothing was omitted from any retained span, so the package carries no omission record.  No retained line contains a citation, so the wrapper carries no bibliography and no bibliography entry.  No retained line contains a figure environment, an image-inclusion command or drawing code, so no image is delivered and the package declares no asset dependency.  Editorial formatting only, in addition: an AMS \texttt{amsart} preamble that restates the article's own declarations and macros used by the retained lines, namely the environments \texttt{prop} on separate counters, together with the standard packages those lines need.  The retained environments print in this document as Proposition 1 in order of appearance, whereas the article prints the same items with the numbering of its own full document.  The complete native source of the article as distributed in the arXiv e-print for this exact version is bundled unchanged as \texttt{sources/197456/source0.tex}.
\begin{prop}
The generating function $M$ satisfies the following two equations:
\begin{equation}
    c(y - Y)(x\widetilde{S}_+ - \widetilde{R}) = M\left(cx(\widetilde{V}'+cx-f_{10})+\widetilde{R}\omega\right) - c^2xf_{00} + x\widehat{P} + \widetilde{R}(-V'-cy+f_{01}),
\label{eq:first-equation-in-M}
\end{equation}
\begin{equation}
    M(\widehat{f}\omega + cxY) = -xR + \widehat{f}(W-f_{01}).
    \label{eq:second-equation-in-M}
\end{equation}
\end{prop}

\begin{proof}
Let us first write
\begin{align*}
    &\brak{\frac{1}{x-A}\frac{1}{\omega - BA}B\frac{\widetilde{V}'(y)-\widetilde{V}'(B)}{y-B}}= \brak{\frac{1}{x-A}\frac{1}{\omega - BA}(\widetilde{V}'(y)-\widetilde{V}'(B))\left(\frac{y}{y-B}-1\right)} \\
    &= y\brak{\frac{1}{x-A}\frac{1}{\omega - BA}\frac{\widetilde{V}'(y)-\widetilde{V}'(B)}{y-B}} - \brak{\frac{1}{x-A}\frac{1}{\omega - BA}(\widetilde{V}'(y)-\widetilde{V}'(B))} \\
    &= y\widetilde{S}_+ - \widetilde{V}'M + \brak{\frac{1}{x-A}\frac{1}{\omega - BA}\widetilde{V}'(B)} \\
    &=  y\widetilde{S}_+ - \widetilde{V}'M - \left(c\brak{\frac{1}{x-A}\frac{1}{\omega - BA}A} - \brak{\frac{1}{x-A}\frac{1}{\omega - BA}}\brak{\frac{1}{\omega - BA}A} \right)\\
&\gr{\Bigg[\uparrow\textnormal{Splitting of }\brak{\frac{1}{x-A}\frac{1}{\omega - BA}\mathbf{A}}\Bigg]}\\ 
    &= y\widetilde{S}_+ - \widetilde{V}'M - c\brak{\left(\frac{x}{x-A}-1\right)\frac{1}{\omega - BA}} + Mf_{10} = y\widetilde{S}_+ - \widetilde{V}'M - cxM + cf_{00} + Mf_{10}.
\end{align*}

On the other hand:
\begin{align*}
 &c\brak{\frac{1}{x-A}\frac{1}{\omega - BA}B\frac{\widetilde{V}'(y)-\widetilde{V}'(B)}{y-B}} \\
    &=\brak{\frac{1}{x-A}\frac{1}{\omega - BA}V'(A)\frac{\widetilde{V}'(y)-\widetilde{V}'(B)}{y-B}} + \brak{\frac{1}{x-A}}\brak{\frac{1}{x-A}\frac{1}{\omega - BA}\frac{\widetilde{V}'(y)-\widetilde{V}'(B)}{y-B}} \\
    &\hspace{4em}+ \brak{\frac{1}{x-A}\frac{1}{\omega - BA}B}\brak{\frac{1}{\omega - BA}\frac{\widetilde{V}'(y)-\widetilde{V}'(B)}{y-B}} \\
&\gr{\Bigg[\uparrow\textnormal{Splitting of } c\brak{\frac{1}{x-A}\frac{1}{\omega - BA}\mathbf{B}\frac{\widetilde{V}'(y)-\widetilde{V}'(B)}{y-B}}\Bigg]}\\ 
    &= -V'(x)\widetilde{S}_++\widehat{P} + W\widetilde{S}_+ +\brak{\frac{1}{x-A}\frac{1}{\omega - BA}B}\widetilde{R} \\
    &= -V'(x)\widetilde{S}_++\widehat{P}+W\widetilde{S}_++\brak{\frac{1}{x}\left(1+\frac{A}{x-A}\right)\frac{1}{\omega - BA}B}\widetilde{R} \\
    &= -V'(x)\widetilde{S}_++\widehat{P}+W\widetilde{S}_++\frac{\widetilde{R}}{x}\left(f_{01} + \brak{\frac{1}{x-A}\frac{1}{\omega - BA}BA}\right) \\
    &= -V'(x)\widetilde{S}_++\widehat{P}+W\widetilde{S}_++\frac{\widetilde{R}}{x}\left(f_{01} + \brak{\frac{1}{x-A}\left(\frac{\omega}{\omega - BA}-1\right)}\right) \\
    &= -V'(x)\widetilde{S}_++\widehat{P}+W\widetilde{S}_++\frac{\widetilde{R}}{x}\left(f_{01} + \omega M - W\right).
\end{align*}

Combining the two previous equations, we obtain:
\begin{align*}
    c\left(y\widetilde{S}_+ - \widetilde{V}'M - cxM + cf_{00} + Mf_{10}\right) = -V'(x)\widetilde{S}_++\widehat{P}+W\widetilde{S}_++\frac{\widetilde{R}}{x}\left(f_{01} + \omega M - W\right),
\end{align*}
which can be rewritten as
\begin{align*}
    \widetilde{S}_+(cy + V' - W) = M\left(c\widetilde{V}'+c^2x-cf_{10}+\frac{\widetilde{R}\omega}{x}\right) - c^2f_{00} + \widehat{P} + \frac{\widetilde{R}}{x}(f_{01}-W),
\end{align*}
which yields equation~\eqref{eq:first-equation-in-M}. 

Let us now turn to~\eqref{eq:second-equation-in-M}. We have:
\begin{align*}
&c\brak{\frac{1}{x-A}\frac{1}{\omega-BA}B} \\
&= -\brak{\frac{1}{x-A}\frac{1}{\omega-BA}V'(A)} + \brak{\frac{1}{x-A}}\brak{\frac{1}{x-A}\frac{1}{\omega-BA}} + \brak{\frac{1}{x-A}\frac{1}{\omega-BA}B}\brak{\frac{1}{\omega-BA}}\\
&\gr{\Bigg[\uparrow\textnormal{Splitting of } c\brak{\frac{1}{x-A}\frac{1}{\omega-BA}\mathbf{B}}\Bigg]}\\
&=- V'M + R + WM +\brak{\frac{1}{x-A}\frac{1}{\omega-BA}B}f_{00},
\end{align*}
so that
\begin{align*}
    \brak{\frac{1}{x-A}\frac{1}{\omega-BA}B} = -\frac{1}{\widehat{f}}(R+cYM).
\end{align*}
On the other hand,
\begin{align*}
    &\brak{\frac{1}{x-A}\frac{1}{\omega-BA}B} = \brak{\frac{1}{x}\left(1+\frac{A}{x-A}\right)\frac{1}{\omega-BA}B} = \frac{1}{x}\left(f_{01} + \brak{\frac{1}{x-A}\frac{1}{\omega-BA}BA}\right) \\
    &= \frac{1}{x}\left(f_{01} + \brak{\frac{1}{x-A}\left(\frac{\omega}{\omega-BA}-1\right)}\right)= \frac{1}{x}\left(f_{01}+\omega M - W\right).
\end{align*}

Combining the two previous equations, we obtain:
\begin{equation*}
    x(R + cYM) = -\widehat{f}(f_{01}+\omega M - W),
\end{equation*}
which yields~\eqref{eq:second-equation-in-M}.

\end{proof}

\end{document}
